% Options for packages loaded elsewhere
\PassOptionsToPackage{unicode}{hyperref}
\PassOptionsToPackage{hyphens}{url}
\documentclass[
  ukrainian,
]{article}
\usepackage{xcolor}
\usepackage[margin=2.5cm]{geometry}
\usepackage{amsmath,amssymb}
\setcounter{secnumdepth}{-\maxdimen} % remove section numbering
\usepackage{iftex}
\ifPDFTeX
  \usepackage[T1]{fontenc}
  \usepackage[utf8]{inputenc}
  \usepackage{textcomp} % provide euro and other symbols
\else % if luatex or xetex
  \usepackage{unicode-math} % this also loads fontspec
  \defaultfontfeatures{Scale=MatchLowercase}
  \defaultfontfeatures[\rmfamily]{Ligatures=TeX,Scale=1}
\fi
\usepackage{lmodern}
\ifPDFTeX\else
\setmainfont{DejaVu Serif}
\setsansfont{DejaVu Sans}
\setmonofont{DejaVu Sans Mono}
\setmathfont{STIX Two Math}
\fi
\ifPDFTeX\else
  % xetex/luatex font selection
\fi
% Use upquote if available, for straight quotes in verbatim environments
\IfFileExists{upquote.sty}{\usepackage{upquote}}{}
\IfFileExists{microtype.sty}{% use microtype if available
  \usepackage[]{microtype}
  \UseMicrotypeSet[protrusion]{basicmath} % disable protrusion for tt fonts
}{}
\makeatletter
\@ifundefined{KOMAClassName}{% if non-KOMA class
  \IfFileExists{parskip.sty}{%
    \usepackage{parskip}
  }{% else
    \setlength{\parindent}{0pt}
    \setlength{\parskip}{6pt plus 2pt minus 1pt}}
}{% if KOMA class
  \KOMAoptions{parskip=half}}
\makeatother
\usepackage{graphicx}
\graphicspath{{/Users/kalinicenkovitalijmikolajovic/.codex/worktrees/5c15/Script_for_translate_book/Output_Final/}}
\setkeys{Gin}{width=0.74\linewidth,height=0.78\textheight,keepaspectratio}
\makeatletter
\newsavebox\pandoc@box
\newcommand*\pandocbounded[1]{% scales image to fit in text height/width
  \sbox\pandoc@box{#1}%
  \Gscale@div\@tempa{\textheight}{\dimexpr\ht\pandoc@box+\dp\pandoc@box\relax}%
  \Gscale@div\@tempb{\linewidth}{\wd\pandoc@box}%
  \ifdim\@tempb\p@<\@tempa\p@\let\@tempa\@tempb\fi% select the smaller of both
  \ifdim\@tempa\p@<\p@\scalebox{\@tempa}{\usebox\pandoc@box}%
  \else\usebox{\pandoc@box}%
  \fi%
}
% Set default figure placement to htbp
\def\fps@figure{htbp}
\makeatother
\ifLuaTeX
\usepackage[bidi=basic,shorthands=off,provide=*]{babel}
\else
\usepackage[bidi=default,shorthands=off,provide=*]{babel}
\fi
\ifLuaTeX
  \usepackage{selnolig} % disable illegal ligatures
\fi
\setlength{\emergencystretch}{3em} % prevent overfull lines
\providecommand{\tightlist}{%
  \setlength{\itemsep}{0pt}\setlength{\parskip}{0pt}}
\usepackage{bookmark}
\IfFileExists{xurl.sty}{\usepackage{xurl}}{} % add URL line breaks if available
\urlstyle{same}
\hypersetup{
  pdflang={uk},
  hidelinks,
  pdfcreator={LaTeX via pandoc}}

\author{}
\date{}

\begin{document}

\pandocbounded{\includegraphics[keepaspectratio]{images/_page_205_Figure_5.jpeg}}

Фіг. 28 1. Типова залежність потоку від густини.

Тепер ми готові застосувати \texttt{u\_i+f\_x=0} рівняння до проблеми
дорожнього руху. Згідно з правилом диференціації комплексної функції
\texttt{f\_x=\textbackslash{}frac\{df\}\{du\}\textbackslash{},u\_x} ; Це
означає, що закон збереження можна переписати у вигляді

\begin{verbatim}
u_t + \frac{df}{du} u_x = u_t + g(u) u_x = 0.
\end{verbatim}

Пусть, например, зависимость потока от плотности квадратичная, т. е.
\texttt{f(u)\ =\ u\^{}\{2\}} , тогда закон сохранения принимает
следующую форму:

\begin{verbatim}
u_t + 2uu_x = 0.
\end{verbatim}

Якщо початкова густина розподілу вагонів була
\texttt{u(x,\ 0)\ =\ \textbackslash{}varphi(x)} , то для знаходження
густини розподілу в довільний момент часу необхідно розв'язати задачу з
початковою умовою (задача Коші):

\begin{verbatim}
\begin{array}{llll} (\text{УЧП}) & u_t + 2uu_x = 0, & -\infty < x < \infty, & 0 < t < \infty, \\ (\text{HV}) & u\left(x, \ 0\right) = \varphi\left(x\right), & -\infty < x < \infty. \end{array}
\end{verbatim}

Тепер розглянемо розв'язок нелінійної задачі Коші.

\section{Нелінійна задача з початковою умовою (задача
Коші)}\label{ux43dux435ux43bux456ux43dux456ux439ux43dux430-ux437ux430ux434ux430ux447ux430-ux437-ux43fux43eux447ux430ux442ux43aux43eux432ux43eux44e-ux443ux43cux43eux432ux43eux44e-ux437ux430ux434ux430ux447ux430-ux43aux43eux448ux456}

Рассмотрим задачу

\begin{verbatim}
\begin{array}{llll} (\mathsf{Y}\mathsf{P}\Pi) & u_t + g\left(u\right)u_x = 0, & -\infty < x < \infty, & 0 < t < \infty, \\ (\mathsf{H}\mathsf{Y}) & u\left(x, \ 0\right) = \varphi\left(x\right), & -\infty < x < \infty. \end{array}
\end{verbatim}

Вспомним предыдущую лекцию, где мы изучали уравнение конвективного
переноса:

\begin{verbatim}
u_t + vu_x = 0
\end{verbatim}

.

У цьому рівнянні \texttt{u(x,\ t)} --- це концентрація матерії в потоку,
що рухається зі швидкістю v

\begin{verbatim}
u_1 + g(u)u_x = 0
\end{verbatim}

та рівняння транспорту. Припустимо, що частина води в \texttt{x\_0}
рухається зі швидкістю \texttt{g(u)} (вздовж потоку або проти нього).
Тоді за t секунд координата частини буде

\begin{verbatim}
x = x_0 + g(u) t
\end{verbatim}

(уравнение характеристик).

Пам'ятайте, що концентрація \texttt{u\ (x,\ t)} не змінюється вздовж
характеристики. Якщо ми знаємо початкову концентрацію
\texttt{u(x\_a,\ 0)} ,

\pandocbounded{\includegraphics[keepaspectratio]{images/_page_206_Figure_8.jpeg}}

Фіг. 28.2. Характеристики рівняння \texttt{u\_t+g} (u) \texttt{u\_x=0} ;
a --- загальна форма характеристики, що випливає з характеристики (
\texttt{x\_0} , 0), вздовж характеристики розв'язок залишається сталим і
рівним значенню \texttt{u\ (} x\_0 \texttt{,\ 0)}; b --- приклад
характеристики, що починається на t. (2, 0), рівняння
\texttt{u\_t+3uu\_x=0} з \texttt{u\ (2,\ 0)}=1.

то уравнение характеристик принимает вид
\texttt{x\ =\ x\_0\ +\ g{[}u(x\_0,\ 0){]}\ t} (уравнение характеристики,
начинающейся в точке \texttt{(x\_0,\ 0)} ).

Пусть, например, мы решаем задачу

(УЧП)

\begin{verbatim}
u_t + 3uu_x = 0
\end{verbatim}

,\\
(НУ) u(x, 0) =

MATHBLK0008X

Щоб знайти характеристику, що виходить із точки, наприклад (2, 0),
потрібно записати

\begin{verbatim}
x = 2 + g[u(2, 0)]t =
\end{verbatim}

= \texttt{2\ +\ g{[}1{]}t\ =} = \texttt{2\ +\ 3t} .

Можна стверджувати, що розв'язок нашої задачі \texttt{u(x,\ t)} дорівнює
одному на лінін x=2+3t, тобто \texttt{u(2+3t,\ t)}=1 для всіх
t\textgreater0. Фіг. 28.2 показує характеристики загального випадку та
цього конкретного завдання.

\texttt{Тепер\ зрозуміло,\ що,\ знаючи\ характеристики\ в\ кожній\ точці\ і\ знаючи,\ що\ розв\textquotesingle{}язок\ за\ характеристиками\ не\ змінюється,\ можна\ отримати\ розв\textquotesingle{}язок\ \textquotesingle{}u\textbackslash{}left(x,\textbackslash{}\ t\textbackslash{}right)\textquotesingle{}\ у\ будь-який\ момент\ часу\ t.\ Ми\ не\ отримали\ явного\ виразу\ для\ \textquotesingle{}u\textbackslash{}left(x,\textbackslash{}\ t\textbackslash{}right)\textquotesingle{}\ через\ змінні\ x\ і\ t,\ але\ знання\ характеристик\ рівняння\ дозволяє\ розв\textquotesingle{}язати\ деякі\ цікаві\ задачі.}

\section{Задача дорожного
движения}\label{ux437ux430ux434ux430ux447ux430-ux434ux43eux440ux43eux436ux43dux43eux433ux43e-ux434ux432ux438ux436ux435ux43dux438ux44f}

Розглянемо модель руху, у якій потік обчислюється за формулою
\texttt{f(u)=u\^{}2} , і початковою густиною iso-P і C

\pandocbounded{\includegraphics[keepaspectratio]{images/_page_207_Figure_8.jpeg}}

. 28.3. Начальная плотность автомобилей движущихся слева направо.

виведений на Рис. 28.3. Іншими словами, нам потрібно розв'язати проблему

\begin{verbatim}
(\text{YUII}) \quad u_t + 2uu_x = 0, \quad -\infty < x < \infty, \quad 0 < t < \infty,
(\text{HV}) \quad u(x, \ 0) =
\begin{cases} 1, & x \leq 0, \\ 1 - x, & 0 < x < 1, \\ 0, & 1 \leq x. \end{cases}
\end{verbatim}

Почнемо з визначення характеристик, які виникають з \texttt{(x\_0,\ 0)}
точок зору. За \texttt{x\_0\ \textless{}\ 0} ми отримуємо

\begin{verbatim}
x = x_0 + g [u(x_0, 0)] t =
\end{verbatim}

= \texttt{x\_0\ +\ g\ {[}1{]}\ t\ =} = \texttt{x\_0\ +\ 2t} .

Разрешая относительно t, находим уравнение характеристик

\begin{verbatim}
t = \frac{1}{2} (x - x_0).
\end{verbatim}

Ці прямі лінії показані на рис. 28.4. Тепер розглянемо точки
\texttt{0\ \textless{}\ x\_0\ \textless{}\ 1} . Тут характеристики
визначаються такими співвідношеннями:

\begin{verbatim}
x = x_0 + g[u(x_0, 0)] t =
= x_0 + g[(1 - x_0)] t =
= x_0 + 2(1 - x_0) t.
\end{verbatim}

Разремая последнее соотношение относительно t, получаем

\begin{verbatim}
t = \frac{x - x_0}{2(1 - x_0)}.
\end{verbatim}

При \texttt{1\ ≤\ x\_0\ \textless{}\ ∞} характеристики

\begin{verbatim}
x = x_0 + g[u(x_0, 0)]t =
\end{verbatim}

= \texttt{x\_0\ +\ g{[}0{]}t\ =} = \texttt{x\_0} призводить до
вертикальних ліній, що виходять з xo. Усі характеристики задачі показані
на рис. 28.4.

\pandocbounded{\includegraphics[keepaspectratio]{images/_page_208_Figure_10.jpeg}}

Фіг. 28.4. Характеристики рівняння \texttt{u\_t\ +\ 2uu\_x\ =\ 0} .

Ситуація на дорозі при 0 \textless{} t \textless{} 1/2 досить зрозуміла.
Рис. 28.5 показує розв'язок нашої задачі в різні моменти часу. Зверніть
увагу, що характеристики сходяться в одній точці t=1/2. Отже, щоб
розв'язати задачу на t\textgreater1/2, потрібно використовувати інший
метод. Коли характеристики сходяться в одній точці, говоримо про ударну
хвилю (розрив розв'язку). Тепер потрібно відповісти на питання: на якій
швидкості є передній край

Ударна хвиля рухатиметься вздовж шосе? Хоча це не очевидно, виявляється,
що швидкість поширення розриву відповідає формулі

\begin{verbatim}
S = \frac{f(u_R) - f(u_L)}{u_R - u_L},
\end{verbatim}

де \texttt{u\_R} і \texttt{u\_L} --- значення розв'язку праворуч і
ліворуч від фронту хвилі, а \texttt{f(u\_R)} і \texttt{f(u\_L)} ---
значення потоку при цих густинах.

\pandocbounded{\includegraphics[keepaspectratio]{images/_page_209_Figure_5.jpeg}}

Фіг. 28.5. Щільність руху в різний час.

У нашому прикладі швидкість поширення щілини виявляється рівною одиниці.
Насправді, якщо згадати цей \texttt{f(u)\ =\ u\^{}2} , то одразу
отримуємо

\begin{verbatim}
S = \frac{0-1}{0-1} = 1
\end{verbatim}

.

Це означає, що при t\textgreater1/2 фронт хвилі рухатиметься зліва
направо зі швидкістю одиниці.

Полное решение задачи

\begin{verbatim}
(\text{Y} \, \Pi) \quad u_t + 2uu_x = 0, \quad -\infty < x < \infty,
(\text{HY}) \quad u(x, \ 0) =
\begin{cases} 1, & x \leq 0, \\ 1 - x, & 0 < x < 1, \\ 0, & 1 < x, \end{cases}
\end{verbatim}

изображено на рис. 28.5. \# ЗАУВАЖЕННЯ 1. Ударна хвиля в нашому прикладі
виникає через те, що потік швидко зростає зі збільшенням густини u. Якщо
потік заданий формулою \texttt{f(u)} = u, то рівняння було б у вигляді
\texttt{u\_1+u\_x=0} , і розв'язок цього рівняння представляє хвилю, що
рухається вправо без спотворень. Якщо на мить замислитися над значенням
формули \texttt{f(u)} = u, стає очевидно, що розв'язок рухатиметься
таким чином.

\begin{enumerate}
\def\labelenumi{\arabic{enumi}.}
\setcounter{enumi}{1}
\tightlist
\item
  Прямой подстановкой решения в виде
\end{enumerate}

\begin{verbatim}
u = \varphi(x - g(u) t)
\end{verbatim}

Ви можете переконатися, що ця формула неявно визначає розв'язок проблеми

\begin{verbatim}
\begin{array}{lll} u_t + g\left(u\right) u_x = 0, & -\infty < x < \infty, & 0 < t < \infty, \\ u\left(x, \ 0\right) = \varphi\left(x\right), & -\infty < x < \infty. \end{array}
\end{verbatim}

Например, решение задачи Коши

\begin{verbatim}
(Y\Pi\Pi) \quad u_t + uu_x = 0,
(HY) \quad u(x, 0) = x
\end{verbatim}

неявно вадается соотношением

\begin{verbatim}
u = \varphi [x - g(u) t] =
= x - g(u) t = x - ut.
\end{verbatim}

У нашому конкретному випадку можливо отримати рішення у явній формі (у
багатьох інших випадках це неможливо)

\begin{verbatim}
u(x, t) = \frac{x}{1+t}.
\end{verbatim}

Переконайтеся, що ця функція справді є вирішенням нашої проблеми.

\section{\texorpdfstring{\textbf{ЗАВДАННЯ}}{ЗАВДАННЯ}}\label{ux437ux430ux432ux434ux430ux43dux43dux44f}

\begin{enumerate}
\def\labelenumi{\arabic{enumi}.}
\tightlist
\item
  Решите задачу Коши
\end{enumerate}

\begin{verbatim}
\begin{array}{llll} (\text{YUII}) & u_t + uu_x = 0, & -\infty < x < \infty, & 0 < t < \infty, \\ (\text{HV}) & u(x, \ 0) = \begin{cases} 0, & x < 0, \\ x, & 0 \leqslant x. \end{cases} \end{array}
\end{verbatim}

Побудуйте хронологію рішення для різних моментів часу. Як ви тлумачите
це рішення? Як виглядає зв'язок між потоком і густиною в цій задачі?

\begin{enumerate}
\def\labelenumi{\arabic{enumi}.}
\setcounter{enumi}{1}
\tightlist
\item
  Решите задачу
\end{enumerate}

\begin{verbatim}
\begin{split} & (\text{УЧП}) \quad u_t + u^2 u_x = 0, \quad -\infty < x < \infty, \quad 0 < t < \infty, \\ & (\text{НУ}) \quad u_t(x, 0) = \begin{cases} 0, & x < 0, \\ x, & 0 \leqslant x. \end{cases} \end{split}
\end{verbatim}

Який зв'язок між потоком і щільністю? Чи співпадає поведінка отриманого
вами рішення з тим, чого ви очікували? Порівняйте рішення задачі 1 і
завдання 2.

\begin{enumerate}
\def\labelenumi{\arabic{enumi}.}
\setcounter{enumi}{2}
\tightlist
\item
  Припустимо, що нев'язка рідина протікає через трубу і одночасно
  просочується крізь стінки згідно із законом \texttt{F(u)} = ku
  (г/см·с). Закон збереження (або, точніше, закон зміни) у цьому випадку
  записується як
\end{enumerate}

\begin{verbatim}
u_t + f_x = -F(u)
\end{verbatim}

.

Знайдіть розв'язок цього рівняння, якщо \texttt{f(u)} = u і
\texttt{\textbackslash{}varphi(x)\ =\ u} (x, 0) (початковий розподіл
задається загальним чином). Дайте тлумачення отриманого рішення.

\begin{enumerate}
\def\labelenumi{\arabic{enumi}.}
\setcounter{enumi}{3}
\item
  Знайти розв'язок попереднього рівняння, якщо втрати визначаються
  функцією \texttt{F(x,\ t)} = 1/x. Перевірте розв'язок. Ви розумієте
  його фізичне значення?
\item
  Перевірити, що функція та неявно задана відношенням
\end{enumerate}

\begin{verbatim}
u = \varphi[x - g(u) t],
\end{verbatim}

є розв'язком нелінійної задачі

\begin{verbatim}
\begin{aligned}
(\text{УЧП})\quad & u_t + g(u)u_x = 0, \\
(\text{НУ})\quad & u(x,0) = \varphi(x).
\end{aligned}
\end{verbatim}

\section{СИСТЕМИ РІВНЯНЬ З ЧАСТКОВИМИ
ПОХІДНИМИ}\label{ux441ux438ux441ux442ux435ux43cux438-ux440ux456ux432ux43dux44fux43dux44c-ux437-ux447ux430ux441ux442ux43aux43eux432ux438ux43cux438-ux43fux43eux445ux456ux434ux43dux438ux43cux438}

Мета лекції - показати, що багато фізичних систем, особливо в
гідродинаміці, не можна описати одним рівнянням. Для цього потрібні
системи пов'язаних рівнянь для тиску, густини, температури та компонент
швидкості.

Лекція покаже, як можна розв'язати лінійну систему рівнянь

\begin{verbatim}
u_t + Au_x = 0
\end{verbatim}

якщо перетворити її на нову систему незалежних рівнянь

\begin{verbatim}
v_1 + \Lambda v_x = 0,
\end{verbatim}

а затем независимо решить каждое уравнение этой системы.

У багатьох галузях науки існують системи величин, які взаємопов'язані не
одним, а кількома відношеннями одночасно. Наприклад, у гідродинамікі це
чотири рівняння

\begin{verbatim}
u_t + uu_x + vu_y + \frac{1}{\rho} p_x = 0
\end{verbatim}

\texttt{(збереження\ імпульсу\ вздовж\ осі\ \textquotesingle{}x\textquotesingle{}\ ),\ (29,1)\ \textquotesingle{}v\_t\ +\ uv\_x\ +\ vv\_y\ +\ \textbackslash{}frac\{1\}\{\textbackslash{}rho\}\ p\_y\ =\ 0\textquotesingle{}\ (збереження\ імпульсу\ вздовж\ осі\ \textquotesingle{}y\textquotesingle{}\ ),\ (збереження\ маси),\ (збереження\ маси),\ (збереження\ маси),\ \textquotesingle{}\textbackslash{}left(\textbackslash{}frac\{p\}\{\textbackslash{}rho\^{}\{\textbackslash{}gamma\}\}\textbackslash{}right)\_t\ +\ u\textbackslash{}left(\textbackslash{}frac\{p\}\{\textbackslash{}rho\^{}\{\textbackslash{}gamma\}\}\textbackslash{}right)\_x\ +\ v\textbackslash{}left(\textbackslash{}frac\{p\}\{\textbackslash{}rho\^{}\{\textbackslash{}gamma\}\}\textbackslash{}right)\_y\ =\ 0\textquotesingle{}\ (збереження\ енергії),\ \textquotesingle{}p(x,\ y,\ z,\ t)\textquotesingle{}\ —\ тиск\ у\ рідині,\ \textquotesingle{}u(x,\ y,\ z,\ t)\textquotesingle{}\ —\ \textquotesingle{}x\textquotesingle{}\ компонента\ швидкості,\ \textquotesingle{}v(x,\ y,\ z,\ t)\textquotesingle{}\ —\ \textquotesingle{}y\textquotesingle{}\ компонента\ швидкості,\ \textquotesingle{}\textbackslash{}rho(x,\ y,\ z,\ t)\textquotesingle{}\ —\ густина\ рідини.}

Задана нелінійна система рівнянь називається рівняннями Ейлера руху
рідини. Завдання полягає в одночасному знаходженні всіх невідомих
функцій p, i, v і p, які задовольняють усім чотирьом рівнянням (разом із
початковими та граничними умовами).

Існують й інші причини вивчати системи рівнянь. У теорії звичайних
диференціальних рівнянь (якщо читач пам'ятає) показано, що одне рівняння
другого порядку можна представити як систему з двох рівнянь першого
порядку. Хоча теорія диференціальних рівнянь у похідних не така проста,
дуже часто одне рівняння з похідними у похідних порядків зводиться до
системи диференціальних рівнянь першого порядку. Наприклад, телеграфне
рівняння

\begin{verbatim}
u_{tt} = c^2 u_{xx} + a u_t + b u
\end{verbatim}

путем введения трех новых переменных

\begin{verbatim}
u_1 = u_1
u_2 = u_2
u_3 = u_4
\end{verbatim}

можно свести к системе уравнений первого порядка

(29.2)

\begin{verbatim}
\begin{cases} \frac{\partial u_1}{\partial x} = u_2, \\ \frac{\partial a_1}{\partial t} = u_3, \\ \frac{\partial u_3}{\partial t} = c^2 \frac{\partial u_2}{\partial x} + au_3 + bu_4. \end{cases}
\end{verbatim}

Для розв'язання систем типу (29.2) необхідно добре знати такі поняття,
як власні значення та власні вектори матриці. З цієї причини ми
пропонуємо короткий огляд лінійної алгебри.

\section{Огляд лінійної
алгебри}\label{ux43eux433ux43bux44fux434-ux43bux456ux43dux456ux439ux43dux43eux457-ux430ux43bux433ux435ux431ux440ux438}

Общий случай

Частный случай

Визначення. Якщо A --- квадратна матриця порядку n, то власні значення
матриці називаються n коренями з

Пусть

\begin{verbatim}
A =
\begin{bmatrix} 1 & 1 \\ 4 & 1 \end{bmatrix}
\end{verbatim}

.

ристического уравнения

\begin{verbatim}
\det\left(A-\lambda I\right)=0,
\end{verbatim}

\begin{verbatim}
где  \det(A-\lambda I)  — определитель матрицы  A-\lambda I  .
\end{verbatim}

Деякі власні значення можуть збігатися.

Тогда собственные значения являются корнями уравнения

\begin{verbatim}
\det
\begin{bmatrix} 1-\lambda & 1\\ 4 & 1-\lambda \end{bmatrix}
= (1-\lambda)^2 - 4 =
= \lambda^2 - 2\lambda - 3 = 0.
\end{verbatim}

Откуда

\begin{verbatim}
\lambda_1 = -1
\end{verbatim}

\texttt{,\ \textquotesingle{}\textbackslash{}lambda\_2\ =\ 3\textquotesingle{}\ .}

Визначення. Якщо \texttt{\textbackslash{}lambda} --- власне значення
матриці A, то відповідний власний вектор є таким ненульовим вектором, що
задовольняє відношення

\begin{verbatim}
Ax = \lambda x
\end{verbatim}

.

Приклад. Вектор X =

MATHBLK0047X

\texttt{є\ власним\ вектором,\ що\ відповідає\ власному\ значенню}
\lambda\_2 = 3 , оскільки

\begin{verbatim}
\begin{bmatrix} 1 & 1 \\ 4 & 1 \end{bmatrix}
\begin{bmatrix} 1 \\ 2 \end{bmatrix}
= 3
\begin{bmatrix} 1 \\ 2 \end{bmatrix}
\end{verbatim}

.

Определение. Матрица \texttt{A\^{}\{-1\}} называется обратной к матрице
A, если

\begin{verbatim}
AA^{-1} = A^{-1}A = I_1
\end{verbatim}

Пример, Обратной к

\begin{verbatim}
A =
\begin{bmatrix} 2 & -2 \\ 1 & 1 \end{bmatrix}
\end{verbatim}

является матрица

\begin{verbatim}
A^{-1} =
\begin{bmatrix} 0.25 & 0.5 \\ -0.25 & 0.5 \end{bmatrix}
\end{verbatim}

,

де I - одинична матриця.

поскольку

\begin{verbatim}
AA^{-1} = A^{-1}A =
\begin{bmatrix} 1 & 0 \\ 0 & 1 \end{bmatrix}
\end{verbatim}

.

\section{Діагональна
матриця}\label{ux434ux456ux430ux433ux43eux43dux430ux43bux44cux43dux430-ux43cux430ux442ux440ux438ux446ux44f}

\begin{verbatim}
Нехай A — квадратна матриця порядку n і нехай всі її власні значення  \lambda_1  ,  \lambda_2  ...  \lambda_n  різні. Побудуємо матрицю P так, щоб у неї були *стовпці* координати власних векторів (k-й стовпець містить власний вектор  X_k  , що відповідає власному значенню  \lambda_k  ), тобто
\end{verbatim}

\begin{verbatim}
P =
\begin{bmatrix} X_1 & X_{\tilde{2}} & \dots & X_n \\ X_{1} & X_{2} & \dots & X_n \end{bmatrix}
=
\begin{bmatrix} x_{11} & x_{1\tilde{2}} & \dots & x_{1n} \\ x_{21} & x_{22} & \dots & x_{2n} \\ \vdots & \vdots & \ddots & \vdots \\ x_{n1} & x_{n\tilde{3}} & \dots & x_{nn} \end{bmatrix}
\end{verbatim}

.

Тогда матрица Л, определяемая соотношением

\begin{verbatim}
\Lambda = P^{-1}AP,
\end{verbatim}

где \texttt{P\^{}\{-1\}} --- обратная к матрице P, будет диагональной
матрицей вида

\begin{verbatim}
\Lambda =
\begin{bmatrix} \lambda_1 & & & \\ & \lambda_2 & 0 & \\ & & \ddots & \\ & & & \lambda_n \end{bmatrix}
\end{verbatim}

.

Например, у матрицы

\begin{verbatim}
A =
\begin{bmatrix} 0 & 8 \\ 2 & 0 \end{bmatrix}
\end{verbatim}

власні значення дорівнюють \texttt{\textbackslash{}lambda\_\ 1\ =\ 4} і
\texttt{\textbackslash{}lambda\_2\ =\ -4} , а власні вектори дорівнюють
та власним векторам відповідно

\begin{verbatim}
X_1 =
\begin{bmatrix} 2 \\ 1 \end{bmatrix}
in  X_2 =
\end{verbatim}

\pandocbounded{\includegraphics[keepaspectratio]{images/_page_221_Picture_9.jpeg}}

Рис. 30.1. Қолеблющаяся мембрана (гиперболическая смешанная задача).

где

\begin{verbatim}
\Delta U = U_{rr} + \frac{1}{r} U_r + \frac{1}{r^2} U_{\theta\theta}.
\end{verbatim}

\begin{verbatim}
Зверніть увагу, що тут ми вимагаємо, щоб константа розділення була додатною (тому позначимо її  \lambda^{2-1}  )). Лише в цьому випадку функції T(t) будуть періодичними.
\end{verbatim}

\begin{verbatim}
На следующем шаге мы решим уравнение Гельмгольца, но сначала разберемся с граничным условием для него. Чтобы найти его, подставим  u(r,  \theta, t) = U(r, \theta) T(t)  в граничное условие для мембраны
\end{verbatim}

\begin{verbatim}
u(1, \theta, t) = U(1, \theta) T(t) = 0, \quad 0 < t < \infty,
\end{verbatim}

или

\begin{verbatim}
U(1, \theta) = 0.
\end{verbatim}

Следовательно, для того чтобы найти форму фундаментальных колебаний
мембраны, необходимо решить задачу

\begin{verbatim}
\Delta U + \lambda^2 U = 0,
U(1, \theta) = 0.
\begin{cases} \Delta U + \mu u = 0 & \text{в круге,} \\ U = 0 & \text{на границе круга} \end{cases}
\end{verbatim}

\begin{verbatim}
при  \mu \le 0  имеет только тривиальное решение U=0. — Прим. ред.
\end{verbatim}

1) Поскольку краевая задача

\begin{verbatim}
Это хорошо известная эллиптическая задача на собственные значения, и нам надо найти все  \lambda , при которых она имеет ненуленые решения. Решения  U\left(r,\,\theta\right)  описывают форму фундаментальных колебаний мембраны, а собственные числа  \lambda  являются нулями некоторых функций Бесселя и пропорциональны частотам этих колебаний.
\end{verbatim}

Итак, нам предстоит теперь решить задачу на собственные значения для
уравнения Гельмгольца (эта задача очень важна и сама по себе). Будем
решать ее тем же методом, что и другие линейные однородные уравнения с
нулевыми граничными условиями, а именно методом разделения переменных.

\section{Решение задачи на собственные значения для уравнения
Гельмгольца}\label{ux440ux435ux448ux435ux43dux438ux435-ux437ux430ux434ux430ux447ux438-ux43dux430-ux441ux43eux431ux441ux442ux432ux435ux43dux43dux44bux435-ux437ux43dux430ux447ux435ux43dux438ux44f-ux434ux43bux44f-ux443ux440ux430ux432ux43dux435ux43dux438ux44f-ux433ux435ux43bux44cux43cux433ux43eux43bux44cux446ux430}

Чтобы решить задачу

\begin{verbatim}
\begin{array}{ll} (\mathsf{Y}\mathsf{H}\mathsf{\Pi}) & \Delta U + \lambda^2 U = 0, \\ (\mathsf{\Gamma}\mathsf{Y}) & U(1, \theta) = 0, \end{array}
\end{verbatim}

подставим

\begin{verbatim}
U(r, \theta) = R(r) \Theta(\theta)
\end{verbatim}

в задачу Гельмгольца. Сделав это, получаем

\begin{verbatim}
r^2R'' + rR' + (\lambda^2r^2 - n^2)R = 0
\end{verbatim}

\begin{verbatim}
(уравнение Бесселя),  R(1) = 0  ,  \Theta'' + n^2\Theta = 0  (Зроби все сам). Зверніть увагу, що ми обрали константу розбиття  n^2  ,  n=0,. \ 1,\ 2,\ \ldots  тому, що хочемо, щоб функції  \Theta\left(\theta\right)  були періодичними з періодом  2\pi  . Очевидно, що форма мембрани є періодичною функцією  \theta  . Отже, щоб розв'язати задачу Гельмгольца, потрібно розв'язати два звичайних диференціальні рівняння
\end{verbatim}

\begin{verbatim}
r^2R'' + rR' + (\lambda^2r^2 - n^2) R = 0
\end{verbatim}

\begin{verbatim}
,  0 < r < 1  ,  R(0) < \infty  (физическое ограничение),  R(1) = 0  ,  \Theta'' + n^2\Theta = 0  .
\end{verbatim}

\section{Уравнение
Бесселя}\label{ux443ux440ux430ux432ux43dux435ux43dux438ux435-ux431ux435ux441ux441ux435ux43bux44f}

Уравнение

\begin{verbatim}
r^2R'' + rR' + (\lambda^2r^2 - n^2)R = 0
\end{verbatim}

жорошо известно в теории обыкновенных дифференциальных уравнений. Это
так называемое уравнение Бесселя. Оно имеет два

\begin{verbatim}
линейно-независимых решения  R_{1}(r) = A J_{n}(\lambda r)  функция Бесселя *первого* рода *п*-го порядка,  R_{2}(r) = BY_{n}(\lambda r)  функция Бесселя второго рода n-го порядка,
\end{verbatim}

и, следовательно, общее решение этого уравнения представимо в виде

\begin{verbatim}
R(r) = AJ_n(\lambda r) + BY_n(\lambda r).
\end{verbatim}

\begin{verbatim}
Зверніть увагу, що рішення залежать від параметрів n і  \lambda  , включених у рівняння. Деякі графіки функції Бесселя показані на рис. 30.2. Щоб визначити функції  J_ n(\lambda r)  і  Y_n(\lambda r)  , можливо
\end{verbatim}

\pandocbounded{\includegraphics[keepaspectratio]{images/_page_223_Figure_8.jpeg}}

\pandocbounded{\includegraphics[keepaspectratio]{images/_page_223_Figure_9.jpeg}}

\begin{verbatim}
Рис. 30 2 Функции Бесселя. a — функции Бесселя первого рода;  \pmb{\delta}  — функции Бесселя второго рода.
\end{verbatim}

\begin{verbatim}
використовуйте метод Фробеніуса, який дозволяє знаходити R(r) у вигляді силових рядів. Цей метод містить два лінійно незалежні ряди ступенів  J_n(\lambda r)  та  Y_n(\lambda r)  . Оскільки функція  Y_n(\lambda r)  не обмежена в r=0, ми обираємо розв'язок у вигляді
\end{verbatim}

\begin{verbatim}
R(r) = AJ_n(\lambda r).
\end{verbatim}

\begin{verbatim}
Последний шаг в определении  R\left( r\right)  связан с использованием граничного условия  R\left( \mathbf{l}\right) =0  для определения всех  \lambda  (в данный
\end{verbatim}

Таблица 30.1

Корни \textbar{} уравнений. \textbar{} J\_n(r) = 0 \textbar{}

\textbar-------\textbar------------\textbar--------------\textbar{}

\textbar-------\textbar------------\textbar--------------\textbar{}

Γ \textbar{} 1 \textbar{} n \textbar{} \textbar{} \textbar{} \textbar{}
\textbar{} \textbar{} \textbar{}

\textbar---\textbar---\textbar-------\textbar-------\textbar-------\textbar-------\textbar-------\textbar--\textbar--\textbar{}

~~\textbar{} \textbar{} 0 \textbar{} 1 \textbar{} 2 \textbar{} 3
\textbar{} 4 \textbar{} \textbar{} \textbar{}

~~\textbar{} 1 \textbar{} 2,40 \textbar{} 3,83 \textbar{} 5,13
\textbar{} 6,38 \textbar{} 7,59 \textbar{} \textbar{} \textbar{}

ı \textbar{} 2 \textbar{} 5,52 \textbar{} 7,02 \textbar{} 8,42
\textbar{} 9,76 \textbar{} 11,06 \textbar{} \textbar{} \textbar{}

ľ \textbar{} 3 \textbar{} 8,65 \textbar{} 10,17 \textbar{} 11,62
\textbar{} 13,02 \textbar{} 14,37 \textbar{} \textbar{} \textbar{}

~~\textbar{} 4 \textbar{} 11,79 \textbar{} 13,32 \textbar{} 14,80
\textbar{} 16,22 \textbar{} 17,62 \textbar{} \textbar{} \textbar{}

~~\textbar{} 5 \textbar{} 14,93 \textbar{} 16,47 \textbar{} 17,96
\textbar{} 19,41 \textbar{} 20,83 \textbar{} \textbar{} \textbar{}

~~\textbar{} . \textbar{} • \textbar{} \textbar{} • \textbar{}
\textbar{} • \textbar{} \textbar{} \textbar{}

~~\textbar{} \textbar{} \textbar{} • \textbar{} \textbar{} • \textbar{}
• \textbar{} \textbar{} \textbar{}

~~\textbar{} • \textbar{} • \textbar{} • \textbar{} • \textbar{} •
\textbar{} • \textbar{} \textbar{} \textbar{}

\begin{verbatim}
момент нас не интересует коэффициент A). Подставляя R (1) = 0, в R (r) =  J_n  (  \lambda r  ), получаем  J_n  (  \lambda  ) = 0. Другими словами, для того чтобы функция R (r) обращалась в нуль на границе круга, мы должны выбрать константу разделения так, чтобы она была корнем уравнения  J_n  (r) = 0,  \tau  . e.
\end{verbatim}

\begin{verbatim}
\lambda = k_{nm}
\end{verbatim}

где \texttt{k\_\{nm\}} --- m-й корень уравнения \texttt{J\_n}(r) = 0 .

\begin{verbatim}
Таблиці цих коренів добре відомі, але їх можна обчислити на комп'ютері. Значення перших кількох коренів наведені в Таблиці 1. 30.1. Знаючи ці корені, ми можемо одразу побудувати розв'язок задачі граничного значення для рівняння Гельмгольца. Якщо власні значення  \lambda = k_{nm}  , то власні функції  U_{ nm}(r)  визначаються формулою
\end{verbatim}

\begin{verbatim}
U_{nm}(r, \theta) = J_n(k_{nm}r) [A \sin(n\theta) + B \cos(n\theta)].
\end{verbatim}

\begin{verbatim}
На рис. 30.3 изображены некоторые из этих функций для различных значений n и m. Общая форма функции  U_{nm}(r,\theta)  не зависит от значений констант A и B. Эти константы влияют только на амплитуду колебаний и расположение линии узлов относительно  \theta=0  .
\end{verbatim}

Каждая функция представляет собой фундаментальное колебание круглой
мембраны с частотой

\begin{verbatim}
f_{nm} = k_{nm}c/2\pi.
\end{verbatim}

Эти частоты мы нашли из решений

\begin{verbatim}
T_{nm}(t) = A \sin(k_{nm}ct) + B \cos(k_{nm}ct)
\end{verbatim}

\section{уравнения}\label{ux443ux440ux430ux432ux43dux435ux43dux438ux44f}

\begin{verbatim}
T'' + k_{nm}^2 c^2 T = 0.
\end{verbatim}

\begin{verbatim}
Интересно отметить, что отношение частоты колебания  U_{nm}(\mathbf{r}, \mathbf{\theta})  к частоте основного колебания  U_{nt}(\ma
\end{verbatim}

\begin{verbatim}
sU(x) = -V \frac{dU}{dx}, \quad 0 < x < \infty,
\begin{cases} u(R_1, \theta) = g_1(\theta), \\ u(R_2, \theta) = g_2(\theta). \end{cases}
U(0) = P/s.
\begin{cases} 2/3, & n=1, \\0 & \text{для всіх інших } n, \\ d_n=\begin{cases} -2/3, & n=1, \\0 & \text{для всіх інших} n. \end{cases} \end{cases}
H(\xi) =
\begin{cases} 1, & 0 \le \theta < \pi, \\ 0, & \pi \le \theta < 2\pi. \end{cases}
\begin{cases} 0, & \xi < 0, \\ 1, & \xi \geqslant 0. \end{cases}
\end{verbatim}

\begin{verbatim}
INL1448X (\Gamma Y)   u(1, \varphi) = \cos(3\varphi)  .
\end{verbatim}

\begin{verbatim}
Используя тригонометрию, попытайтесь выразить  \cos{(3\phi)}  через  \cos{\phi}, \, \cos^2{\phi}, \, \cos^3{\phi}, \, a  затем скомбинируйте их так, чтобы по-
\end{verbatim}

лучилось разложение

\begin{verbatim}
\cos(3\varphi) = a_0 P_0(\cos\varphi) + a_1 P_1(\cos\varphi) + \dots
\end{verbatim}

\begin{enumerate}
\def\labelenumi{\arabic{enumi}.}
\setcounter{enumi}{4}
\tightlist
\item
  Решите задачу Дирихле в шаре
\end{enumerate}

\begin{verbatim}
\Delta u = 0, \quad 0 < r < 1, \quad 0 \le \varphi \le \pi, \quad 0 \le \theta < 2\pi.
u(1, \varphi) =
\begin{cases} 1, & 0 \le \varphi \le \pi/2, \\ -1, & \pi/2 < \varphi \le \pi. \end{cases}
\end{verbatim}

В этой задаче верхняя полусфера нагрета (+1), а нижняя охлаждена (-1).

\begin{enumerate}
\def\labelenumi{\arabic{enumi}.}
\setcounter{enumi}{5}
\tightlist
\item
  Найдите решение задачи Дирихле в шаре
\end{enumerate}

(УЧП)

\begin{verbatim}
\Delta u = 0
,  1 < r < \infty  ,  0 \le \varphi \le \pi  ,  0 \le \theta < 2\pi  ,  (\Gamma Y)   u(1, \varphi) = 1 + \cos \varphi  .
\end{verbatim}

Проверьте ответ.

Лекция 36

\section{НЕОДНОРОДНАЯ ЗАДАЧА ДИРИХЛЕ (ФУНКЦИЯ
ГРИНА)}\label{ux43dux435ux43eux434ux43dux43eux440ux43eux434ux43dux430ux44f-ux437ux430ux434ux430ux447ux430-ux434ux438ux440ux438ux445ux43bux435-ux444ux443ux43dux43aux446ux438ux44f-ux433ux440ux438ux43dux430}

МЕТА ЛЕКЦІЇ: Показати, як гетерогенну задачу Діріхле можна розв'язати за
допомогою функції Гріна (вихідної функції). У цьому важливому методі
права частина рівняння розглядається як певна вхідна дія і розкладається
на неперервний набір дельтаподібних джерел, розподілених по заданій
площі. Потім визначається відповідь системи на кожне таке джерело
(функція Гріна), і всі відповіді підсумовуються (інтегровані). У
результаті ми отримуємо повне розв'язання проблеми.

Достаточно общей задачей прикладной математики является определение
потенциала в некоторой области пространства, если задано распределение
источников \texttt{f(x,\ y)} енутри этой области. В электростатике
потенциал (в вольтах) в области D обусловлен распределением зарядов с
плотностью \texttt{f(x,\ y)} по этой области. Типичная задача ---
определение потенциала в круге в двумерном случае.

В такой задаче потенциал должен удовлетворять уравнению

Пуассона (например, с нулевыми граничными условиями)

(36.1) (УЧП)

\begin{verbatim}
u_{rr} + \frac{1}{r} u_r + \frac{1}{r^2} u_{\theta\theta} = f(r, \theta), \quad 0 \le r < 1, \quad 0 \le \theta < 2\pi,
\end{verbatim}

\begin{verbatim}
(ГУ)  u(1, \theta) = 0, \quad 0 \le \theta \le 2\pi.  Ми не випадково обрали нульові умови. Якщо б ми хотіли розв'язати задачу в загальному випадку, з ненульовими граничними умовами та неоднорідним рівнянням, то внесок неоднорідних граничних умов можна було б знайти за допомогою формули інтегралу Пуассона з лекції 33.
\end{verbatim}

Чтобы получить хоть какое-то представление о неоднородных
дифференциальных уравнениях, давайте рассмотрим графически решение
следующей задачи для уравнения Пуассона:

(УЧП)

\begin{verbatim}
\Delta u = -q
\end{verbatim}

\begin{verbatim}
,  0 < r < 1  ,  0 \le \theta \le 2\pi  (  q  —положительная константа),
\end{verbatim}

(HY)

\begin{verbatim}
u(1, \theta) = 0, \quad 0 \le \theta \le 2\pi.
\end{verbatim}

\begin{verbatim}
На межі потенціал (або температура, якщо ви бажаєте) утримується на нулі, а лапласіан дорівнює q у всіх точках області. Оскільки  \Delta u\left(p\right)  слугує мірою різниці між функцією  u\left(p\right)  і її середнім, рівняння Пуассона стверджує, що поверхня  u\left(r\right)  завжди опукла вниз. Інакше кажучи, вона виглядатиме як тонка мембрана з нерухомим краєм, яка здувається зверху вниз потоком повітря.
\end{verbatim}

Якщо права частина рівняння \texttt{f(x,\ y)} змінюється від точки до
точки, то опуклість також буде змінною функцією точки цієї області.

Перейдем теперь к главной части нашей лекции: определению функции Грина
и решению уравнения (36.1).

Сначала, однако, мы должны ввести понятие потенциала точеных источников
и стоков.

\section{Потенциалы точечных источников и
стоков}\label{ux43fux43eux442ux435ux43dux446ux438ux430ux43bux44b-ux442ux43eux447ux435ux447ux43dux44bux445-ux438ux441ux442ux43eux447ux43dux438ux43aux43eux432-ux438-ux441ux442ux43eux43aux43eux432}

Щоб розв'язати неоднорідне лінійне рівняння, достатньо розв'язати це
рівняння за допомогою точкового джерела, оскільки розв'язок задачі в
загальному випадку можна знайти, підсумувавши внески з точкових джерел.
Наше завдання --- визначити потенціал, який створює точкове джерело (або
поглинач) у певній області. Ці джерела можна інтерпретувати по-різному.
У теорії теплопровідності джерело можна розглядати як точку, де походить
тепло, а дренаж --- як точку, де поглинає тепло. В електростатиці
джерело можна розглядати як окремий позитивний заряд (протон), а дренаж
--- як окремий негативний заряд (електрон). У будь-якому разі, незалежно
від інтерпретації, ми знайдемо двовимірний потенціал і,

создаваемый единичным точечным источником (поиск трехмерного потенциала
мы оставляем читателю в качестве упражнения).

Припустимо, що єдине точкове джерело +q розташоване в початку Clear that
heat (або щось інше)

будет вытекать из области вдоль радиальных линий, и, следовательно, для
везничны полного потока, протекающего через окружность радиуса r можно
записать

Полный поток, вытекающий через окруж-

HOCTD

\begin{verbatim}
= -\int_{0}^{2\pi} u_{r}(r) r d\theta = -2\pi r u_{r}(r).
\end{verbatim}

Но полный поток должен совпадать с количеством тепла, созданного внутри
круга (закон сохранения), т. е.

\$\$ \texttt{-2\textbackslash{}pi\ r\ u\_r(r)\ =\ q.} Решая это диф
реренциальное уравнение относительно функции u(r), получаем Р И С

\end{document}
